\documentclass[10pt,a4paper]{amsart}
\usepackage[margin=19mm]{geometry}
\usepackage{amsmath,amssymb,mathtools}
\usepackage[scr=rsfs,cal=euler]{mathalfa}
\usepackage{xfrac}
\usepackage[unicode,hidelinks,bookmarks=false]{hyperref}
\newcommand{\ie}{i.e.\ }
\newcommand{\cf}{cf.\ }
\newcommand{\resp}{respectively\ }
\newcommand{\Subsection}{\subsection}
\providecommand{\nobreakdash}{\nobreak\hbox{-}}
\setlength{\emergencystretch}{2em}
\multlinegap0pt
\usepackage{algorithm}\usepackage{algpseudocode}
\usepackage{amssymb}
\usepackage{bm}
\usepackage{enumerate}
\usepackage{mathtools}
\usepackage{tikz}
\usepackage{xcolor}
\newtheorem{proposition}{Proposition}
\DeclareMathOperator{\Conv}{Conv}
\DeclareMathOperator{\PP}{\mathbb{P}}
\DeclareMathOperator{\RR}{\mathbb{R}}
\DeclareMathOperator{\ZZ}{\mathbb{Z}}
\DeclareMathOperator{\munorm}{\mu_{\text norm}}
\DeclareMathOperator{\muva}{\mu_{\text va}}
\title{Embeddings of weighted projective spaces}
\author{Praise Adeyemo, Dominic Bunnett, Fabi\'an Levic\'an-Santib\'a\~nez}
\date{Source: arXiv:2510.05076v1 (CC BY 4.0)}
\begin{document}
\maketitle

\noindent\emph{Source and licence.} Embeddings of weighted projective spaces. Version arXiv:2510.05076v1 (CC BY 4.0), distributed by arXiv under the Creative Commons Attribution 4.0 International licence, http://creativecommons.org/licenses/by/4.0/, as recorded in the arXiv licence metadata for this exact version and shown on the abstract page https://arxiv.org/abs/2510.05076v1. Full licence notice: \texttt{licenses/arxiv-2510.05076-CC-BY-4.0.txt}.\par\smallskip

\noindent\textit{Editorial scope.} Counted unit: the article's proposition proposition:repeated-weights (source0.tex lines 465--470). The article's complete original proof of it is delivered with it (source0.tex lines 473--496, one \texttt{proof} environment of 24 lines). No mathematics is written, repaired or re-derived: the statement and the proof are the article's own lines, in the article's own notation, and no retained line is substituted or re-flowed.  No further source-local context is needed: every cross-reference of every retained line resolves inside the two delivered spans.  Nothing was omitted from any retained span, so the package carries no omission record.  No retained line contains a citation, so the wrapper carries no bibliography and no bibliography entry.  No retained line contains a figure environment, an image-inclusion command or drawing code, so no image is delivered and the package declares no asset dependency.  Editorial formatting only, in addition: an AMS \texttt{amsart} preamble that restates the article's own declarations and macros used by the retained lines, namely the environments \texttt{proposition} on separate counters, together with the standard packages those lines need.  The retained environments print in this document as Proposition 1 in order of appearance, whereas the article prints the same items with the numbering of its own full document.  The complete native source of the article as distributed in the arXiv e-print for this exact version is bundled unchanged as \texttt{sources/186496/source0.tex}.
\begin{proposition}\label{proposition:repeated-weights}
    Consider the weighted projective spaces $\widetilde{\PP} = \PP(a_0, \dots , a_n)$ and $\PP = \PP(a_0, \dots , a_n , a_n)$, where they only differ by dropping the last repeated weight $a_n$.

    Then 
    \[\muva(\PP) = \muva(\widetilde{\PP}) \quad\text{ and }\quad \munorm(\widetilde{\PP}) = \munorm(\PP) \enspace.\]
\end{proposition}

\begin{proof}
    We translate this into a problem of the corresponding simplices.
    Consider an $n$-tuple of positive integers $b_0, \dots , b_n$ and we denote $\Delta = \Conv(b_0e_{0}, \dots , b_ne_n, b_ne_{n+1}) \subset \RR^{n+2}$ and $\widetilde{\Delta} = \Conv(e_0b_0, \dots , e_nb_n) \subset \RR^{n+1}$.
    Note that $\widetilde{\Delta}$ can be seen as a face of $\Delta$.
    We show that $\Delta$ is very ample if and only if $\widetilde{\Delta}$ is very ample and remark that the proof for normality is the exact same.
    
    We define a map $\widetilde{ \cdot} : \RR^{n+2} \rightarrow \RR^{n+1}$ which sends $(x_0, \dots , x_n,x_{n+1}) \mapsto (x_0 + x_1, \dots , x_n+x_{n+1})$.
    This map restricts to a map $\widetilde{\cdot} : \Delta \rightarrow \widetilde{\Delta}$, mapping lattice points to lattice points.
    Moreover, it surjects onto the lattice points of $\widetilde{\Delta}$ and every fibre over such a lattice point contains at least one lattice point of $\Delta$.
    
    In fact, we can describe the lattice points in the fibres explicitly.
    That is, given $(x_0, \dots , x_n) \in \widetilde{\Delta}\cap \ZZ^{n+1}$, the fibre consists of points
    \[\left\{(x_0, \dots ,x_{n-1},x_{n}-i,i)\,\, | \,\, i=0, \dots , x_n \right\}.\]

    One direction is immediate: $\Delta$ being very ample implies $\widetilde{\Delta}$ is very ample, since it is a face of $\Delta$.
    For the converse, suppose that $\widetilde{\Delta}$ is very ample such that for every $k \geq k_0$ we have the IDP condition.
    Fix such a $k \geq k_0$ and take a lattice point $x \in \Delta \cap \ZZ^{n+1}$.
    Then $\widetilde{x} = \widetilde{x}_1 + \cdots + \widetilde{x}_k$ for $x_i \in \widetilde{\Delta} \cap \ZZ^n$.
    Then we can choose the $x_i$ over each $\widetilde{x}_i$ such that
    $x = x_1 + \cdots + x_k$.
    Thus $\Delta$ is very ample.

    It follows that the normal and very ample indices are the same.
\end{proof}

\end{document}
